%
% 52-meromorph.tex -- Das Argumentprinzip für beliebige meromorphe Funktionen
%
% (c) 2026 Prof Dr Andreas Müller
%
\subsection{Das Argumentprinzip für beliebige meromorphe Funktionen}
%
% fig-argument.tex -- 
%
% (c) 2026 Prof Dr Andreas Müller
%
\begin{figure}
\centering
\includegraphics{chapters/040-analytisch/images/argument.pdf}
\caption{Bild eines Kreises unter einer Funktion $f(z)$ in Abhängigkeit
von der Anzahl der Nullstellen (blaue Kreise) und Pole (blaue Kreuze),
die vom Kreis umlaufen werden.
Die Differenz ist die Anzahl Umläufe der Bildkurve um den Nullpunkt.
\label{buch:analytisch:fig:argument}}
\end{figure}
%
Der Satz~\ref{buch:analytisch:argument:satz:ordnung}
liefert Information für eine einzelne isolierte Singularität.
Wir möchten jetzt zeigen, dass sich die Schlussfolgerung auch
verallgemeinern lässt auf einen Weg, der mehr als eine Null-
oder Polstelle der Funktion $f$ umschliesst.

\begin{satz}[Argumentprintzip für meromorphe Funktionen]
\label{buch:analytisch:argument:satz:argument}
\index{Argumentprinzip!meromorphe Funktionen}%
\index{meromorph}%
Sei jetzt $f$ ein meromorphe Funktion und $\gamma$ ein einfach
geschlossener Weg in $\mathbb{C}$, der keine Nullstellen oder Polstellen
von $f$ trifft.
Dann gilt
\[
\frac{1}{2\pi i}
\oint_\gamma
\frac{f'(z)}{f(z)}
\,dz
=
\sum_{\text{$a$ Singularität von $f$}}
\omega(a;f),
\]
wobei die Summe über alle Null- oder Polstellen $a$ von $f$ 
zu erstrecken ist.
\end{satz}

\begin{proof}
Nach Satz~\ref{buch:integration:homologie:satz:pfadpunkte} ist das
Wegintegral über $\gamma$ gleich gross wie die Summe der Integrale
über kleine, im Gegenuhrzeigersinn durchlaufene Kreise um die einzelnen
Singularitäten.
Die Summanden haben nach
Satz~\ref{buch:analytisch:argument:satz:ordnung}
den Wert $\omega(a;f)$.
Damit ist der Satz bewiesen.
\end{proof}

In Abbildung~\ref{buch:analytisch:fig:argument}
wird das Bild eines Kreises unter einer Funktion $f(z)$ in Abhängigkeit
von der Anzahl der Nullstellen und Pole, die vom Kreis umlaufen werden.
Sie zeigt, dass die Bildkurve den Nullpunkt im Bildbereich soviele
Male umläuft, wie die Differenz der vom Kreis umlaufenen Nullstellen und
Pole angibt.

Der Satz \ref{buch:analytisch:argument:satz:argument} ermöglicht,
die Anzahl der Nullstellen und Polstellen einer Funktion zu zählen,
ohne dass man wissen muss, wo sich diese befinden.
Dies wird zum Beispiel im Algorithmus von Weyl zur numerischen
Bestimmung der komplexen Nullstellen eines Polynoms ausgenutzt, der in
Kapitel~\ref{chapter:weyl} beschrieben wird.


